%; whizzy paragraph -pdf xpdf -latex ./whizzypdfptex.sh
%; whizzy-paragraph "^\\\\begin{frame}\\|\\\\emtext"
% latex beamer presentation.
% platex, latex-beamer でコンパイルすることを想定。 

%     Tokyo Debian Meeting resources
%     Copyright (C) 2012 Junichi Uekawa

%     This program is free software; you can redistribute it and/or modify
%     it under the terms of the GNU General Public License as published by
%     the Free Software Foundation; either version 2 of the License, or
%     (at your option) any later version.

%     This program is distributed in the hope that it will be useful,
%     but WITHOUT ANY WARRANTY; without even the implied warreanty of
%     MERCHANTABILITY or FITNESS FOR A PARTICULAR PURPOSE.  See the
%     GNU General Public License for more details.

%     You should have received a copy of the GNU General Public License
%     along with this program; if not, write to the Free Software
%     Foundation, Inc., 51 Franklin St, Fifth Floor, Boston, MA  02110-1301 USA

\documentclass[cjk,dvipdfmx,12pt]{beamer}
\usetheme{Tokyo}
\usepackage{monthlypresentation}

%  preview (shell-command (concat "evince " (replace-regexp-in-string "tex$" "pdf"(buffer-file-name)) "&")) 
%  presentation (shell-command (concat "xpdf -fullscreen " (replace-regexp-in-string "tex$" "pdf"(buffer-file-name)) "&"))
%  presentation (shell-command (concat "evince " (replace-regexp-in-string "tex$" "pdf"(buffer-file-name)) "&"))

%http://www.naney.org/diki/dk/hyperref.html
%日本語EUC系環境の時
\AtBeginDvi{\special{pdf:tounicode EUC-UCS2}}
%シフトJIS系環境の時
%\AtBeginDvi{\special{pdf:tounicode 90ms-RKSJ-UCS2}}

\newenvironment{commandlinesmall}%
{\VerbatimEnvironment
  \begin{Sbox}\begin{minipage}{1.0\hsize}\begin{fontsize}{8}{8} \begin{BVerbatim}}%
{\end{BVerbatim}\end{fontsize}\end{minipage}\end{Sbox}
  \setlength{\fboxsep}{8pt}
% start on a new paragraph

\vspace{6pt}% skip before
\fcolorbox{dancerdarkblue}{dancerlightblue}{\TheSbox}

\vspace{6pt}% skip after
}
%end of commandlinesmall

\title{東京エリアDebian勉強会}
\subtitle{第110回 2014年2月度 OSC 2014 Tokyo 出張勉強会}
\author{岩松 信洋 / iwamatsu@debian.org}
\date{2014年3月1日}
\logo{\includegraphics[width=8cm]{image200607/openlogo-light.eps}}

\begin{document}

\begin{frame}
\titlepage{}
\end{frame}

\begin{frame}{Agenda}
  \begin{itemize}
   \item Debian Update
   \item Debian の EFI/UEFI 対応について
   \item 質疑応答
   \item おしらせ
  \end{itemize}
\end{frame}
%
%\begin{block}{ブロック}
%\end{block}
%\begin{alertblock}{警告ブロック}
%\end{alertblock}
%\begin{exampleblock}{例ブロック}
%\end{exampleblock}
%
\emtext{Debian Update}

\begin{frame}{Debian Update}

\begin{itemize}
\item ポイントリリース
\begin{itemize}
\item 2014-02-15 Debian 6.0 更新: 6.0.9 リリース
\item 2014-02-08 Debian 7 更新: 7.4 リリース
\item 2013-12-14 Debian 7 更新: 7.3 リリース
\end{itemize}
\end{itemize}
\end{frame}

\begin{frame}{Debian Update}
\begin{itemize}
\item Debian Contributors list \\
\begin{itemize}
\item \url{https://contributors.debian.org/} \\
\item Debian に貢献している人はだれなのか、リストアップ。\\
\item パッケージアップロードや、Wikiなどから算出。\\
\item 残絵なことにパッケージメンテナはリストアップされるようになっていない。
\end{itemize}
\end{itemize}
\end{frame}

\begin{frame}{Debian Update}

\begin{itemize}
\item IA64 should be removed from testing\\
\begin{itemize}
\item 次のリリース(Jessie)から IA64 サポートから外されます。\\
\item ports でメンテナンスされるかも。
\end{itemize}
\end{itemize}

\end{frame}

\begin{frame}{Debian Update}

\begin{itemize}
\item Debian のデフォルト init システムをどうするか\\
\begin{itemize}
\item sysvinit or upstart or openrc or systemd \\
\item Debian tech-ctte とともに議論中\\
\end{itemize}
\end{itemize}

\end{frame}

\begin{frame}{Debian Update}

\begin{itemize}
\item autopkgtest \\
\begin{itemize}
\item 作成されたパッケージを自動的にテストするプロジェクト。\\
\item DEP8(Debian Enhancement Proposals 8)で議論され、現在実験中。
\end{itemize}
\end{itemize}
\end{frame}

\emtext{Debian の EFI/UEFI 対応について}

\begin{frame}[containsverbatim]{Debian の EFI/UEFI 対応について}
\begin{itemize}
\item EFI / EFI の詳細な話はしません。
\item EFI とは？
\begin{itemize}
\item (Unified) Extensible Firmware Interface の略
\item いままで使われてきたBIOSの置き換え
\item 3TB以上のHDDをサポート
\item 様々なデバイスのサポート
\end{itemize}
\end{itemize}

\end{frame}

\begin{frame}[containsverbatim]{Debian の EFI/UEFI 対応について}

\begin{itemize}
\item セキュアブート
\begin{itemize}
\item 利用できるソフトウェア（OSやドライバーなど）を制限する仕組み。
\item ルートキットや不正なドライバーが OS より先に実行されることを防ぐ。
\end{itemize}
\end{itemize}

\end{frame}

\begin{frame}[containsverbatim]{Debian の EFI/UEFI 対応について}

\begin{itemize}
\item Linux のセキュアブート対応
\begin{itemize}
\item Secure Bootをdisableに設定する\\
ベンダー依存でできない場合もあるので注意
\item マイクロソフトの鍵で署名されたブートローダを使用する\\
\item Ubuntu、Fedora は shim と呼ばれるブートローダを使ってブート\\
\begin{enumerate}
\item UEFI
\item shim(マイクロソフトの鍵で署名されたブートローダ)
\item 2nd Boot Loader (grubなど)
\item Linuxカーネル
\item カーネルローダブルモジュール
\end{enumerate}
\item 自分の鍵をUEFIに組み込む
\end{itemize}
\end{itemize}

\end{frame}

\begin{frame}[containsverbatim]{Debian の EFI/UEFI 対応について}

\begin{itemize}
\item Debian のセキュアブート対応
\begin{itemize}
\item 未対応
\item shim が Free ではない
\item 署名機構をどうするか検討中
\item Secure Bootをdisableにしてインストール
\end{itemize}
\end{itemize}

\end{frame}

%\begipn{frame}[containsverbatim]{DebianのEFI対応状況}

%\begin{commandline}
%CONFIG_RELOCATABLE=y
%CONFIG_EFI=y
%CONFIG_EFI_STUB=y
%CONFIG_FB_EFI=y
%CONFIG_FRAMEBUFFER_CONSOLE=y
%\end{commandline}

%\end{frame}

\begin{frame}[containsverbatim]{DebianのEFI対応状況}

\begin{itemize}
\item もちろんデフォルトで対応。しかしスムーズには行かない。\\
インストールするデバイスによって対応が異なるため。（が、インストーラでの対応はできる）
\item インストールぐらいさっさとしたい。\\
EFI がレガシーBIOSを対応している場合がある。$\rightarrow$
レガシーBIOSモードでブートすれば楽ちん！
\item しかしきちんと動かない場合がある
\begin{itemize}
\item EFI（EFIから読み出されるファームウェア）で、デバイスを初期化しているため。
レガシーBIOSではきちんとやってない場合がある。
\end{itemize}
\end{itemize}

\end{frame}

\begin{frame}[containsverbatim]{Debianの対応状況}
\begin{itemize}
\item EFI/UEFI はサポートしている
\item レガシーBIOS動くかもしれないが、きちんと動作しない場合が多い
\item EFIで起動させたい！
\item EFIを使って起動させるには？
\begin{itemize}
\item VAIO Pro の場合
\item Macbook Pro retina の場合
\end{itemize}
\end{itemize}

\end{frame}

\begin{frame}[containsverbatim]{Vaio Proの場合}
\end{frame}

\begin{frame}[containsverbatim]{Vaio Proの場合}

\begin{enumerate}
\item EFIの画面でSecureBoot を外す。
\item Windows8 と Dual Boot したい場合は \texttt{bcdedit /enum all} で ブート マネージャーの確認

\begin{commandline}
$ bcdedit /enum all
（省略）
Windows ブート マネージャー
--------------------------------
identifier              {bootmgr}
device                  partition=\Device\HarddiskVolume3
path                    \EFI\Microsoft\Boot\bootmgfw.efi
description             Windows Boot Manager
locale                  ja-JP
inherit                 {globalsettings}
default                 {current}
resumeobject            {df94d6db-389b-11e3-8f30-9e441e235da0}
displayorder            {current}
toolsdisplayorder       {memdiag}
timeout                 30

(中略)

ファームウェア アプリケーション (101fffff
--------------------------------
identifier              {f652bc01-389b-11e3-be69-806e6f6e6963}
device                  partition=\Device\HarddiskVolume1
path                    \EFI\Microsoft\Boot\bootmgfw.efi
description             Sony Original
\end{commandline}

\end{enumerate}

\end{frame}

\begin{frame}[containsverbatim]{Vaio Proの場合}

インストール時に作成される\texttt{grubx64.efi}を
path（$\backslash$EFI$\backslash$Microsoft$\backslash$Boot$\backslash$bootmgfw.efi）に
コピー。そして再起動。
\begin{commandline}
# mount -t vfat /dev/sda3 /boot/efi
# cd /boot/efi/EFI/Microsoft/Boot
# cp bootmgfw.efi bootmgfw.efi.orig
# cp ../../debian/grubx64.efi bootmgfw.efi
# cd /
# umount /boot/efi
# exit
\end{commandline}

\end{frame}

\begin{frame}[containsverbatim]{Vaio Proの場合}

Debian を再起動し、grub設定にwindows8の設定を追加する。
\begin{commandline}
$ sudo vi /etc/grub.d/40_custom
#!/bin/sh
exec tail -n +3 $0
menuentry "Windows Boot Manager (SONY Original)" {
    insmod part_gpt
    insmod fat
    set root='(hd0,gpt3)'
    search --no-floppy --fs-uuid --set=root 4A3B-7939
    chainloader /EFI/Microsoft/Boot/bootmgfw.efi.orig
}
\end{commandline}
\end{frame}

\begin{frame}[containsverbatim]{Vaio Proの場合}

update-grub を実行する。
\begin{commandline}
$ sudo update-grub
Generating grub.cfg ...
Found background image: /usr/share/images/desktop-base/desktop-grub.png
Found linux image: /boot/vmlinuz-3.10-3-amd64
Found initrd image: /boot/initrd.img-3.10-3-amd64
Found Windows Boot Manager on /dev/sda3@/EFI/Microsoft/Boot/bootmgfw.efi
Adding boot menu entry for EFI firmware configuration
done
\end{commandline}

\end{frame}

\begin{frame}[containsverbatim]{Macbook Pro retina 2013の場合}
\end{frame}

\begin{frame}[containsverbatim]{Macbook Pro retina 2013の場合}

Macbook Proなのに OS X と Dual Bootしない方法。

\begin{itemize}
\item レガシーBIOSでも起動する
\item しかし、SMPで起動しない
\item USBなど、いくつかのデバイスも動作しない
\end{itemize}

\end{frame}

\begin{frame}[containsverbatim]{Macbook Pro retina 2013の場合}

\begin{enumerate}
\item もくもくとインストール。
\item インストーラで GRUB がインストールされた後、「戻る」を選択してシェルに移動
\item chroot 後 /boot/efi/EFI/debian/grubx64.efi
にあるファイルを/boot/efi/EFI/boot/bootx64.efi にコピー

\begin{commandline}
# chroot /target bash
# mkdir -p /boot/efi/EFI/boot
# cp /boot/efi/EFI/debian/grubx64.efi /boot/efi/EFI/boot/bootx64.efi
\end{commandline}

\item chroot とシェルから抜けて、インストーラーを終了する。このときに「option
キー」を押しながら再起動。
\end{enumerate}
\end{frame}

\begin{frame}[containsverbatim]{Macbook Pro retina 2013の場合}

OS X との Dual Bootは？
\begin{itemize}
\item mactel-boot を使う \\
  Intel Mac 専用 EFI ブートローダ\\
  Debianにはまだパッケージがありません。
\item rEFIt、rEFIndを使う \\
  rEFIt は開発終了。引き継いた rEFInd を使いましょう。\\
  DebianにはまだrEFIndパッケージがありません(rEFItはあります)。
\end{itemize}
\end{frame}

\begin{frame}[containsverbatim]{協力}
\begin{itemize}
\item VAIO Pro 11 / dictoss \\
\url{http://pcdennokan.dip.jp/hardware/vaiopro11_debian/}
\item Macbook pro retina \\
\url{https://wiki.debian.org/ja/InstallingDebianOn/Apple/MacBookPro/11-1}
\end{itemize}
\end{frame}

\begin{frame}[containsverbatim]{まとめ}

\begin{itemize}
\item Debian では EFI/UEFI サポートしています
\item Secueboot は未対応
\item レガシーBIOSモードは動かないデバイスが出てくるので注意
\item EFIで動かしたほうがよい
\item ユーザの手を煩わせないように、がんばってインストーラで対応します
\end{itemize}

\end{frame}

\begin{frame}[containsverbatim]{質問}

\begin{center}
なにか質問はありますか？
\end{center}

\end{frame}

\begin{frame}{おしらせ}
\begin{itemize}
\item 出展しています。
\begin{itemize}
\item Wheezy マシン展示
\item ステッカー、チラシ、配布
\item Debianに関するの相談受け付けます。
\item GPGサイン、CAcertサインできます。
\end{itemize}
\item 次回勉強会は 3/15（土）14:00-19:00 スクエアエニックスさん セミナールームで行います。
内容はまだ未定です。
\end{itemize}

\end{frame}

\end{document}

;;; Local Variables: ***
;;; outline-regexp: "\\([ 	]*\\\\\\(documentstyle\\|documentclass\\|emtext\\|section\\|begin{frame}\\)\\*?[ 	]*[[{]\\|[]+\\)" ***
;;; End: ***
